\section{Teknologivalg}
\label{sec:teknologivalg}

Valget av teknologier tok utgangspunkt i kravspesifikasjonen fra HDO
\cite{hdo_kravspec}, som uttrykte et ønske om .NET 10 som rammeverk, samt
krav om containerisering, tokenbasert autentisering og observability.
Gruppen hadde ingen tidligere erfaring med C\# eller .NET, men valgte å
følge HDOs preferanse fordi rammeverket er godt dokumentert og aktivt
vedlikeholdt.

ASP.NET Core Web API ble valgt i controller-basert modus fremfor Minimal
API, da controller-mønsteret gir bedre separasjon av ansvar og er mer
egnet for et API med flere ressurstyper. Keycloak ble valgt som
identitetsleverandør fordi det støtter OAuth 2.0 og JWT Bearer ut av
boksen, i tråd med krav S1.1 \cite{hdo_kravspec}. Swashbuckle genererer
OpenAPI-dokumentasjon direkte fra kildekoden, noe som oppfyller krav D1.1
\cite{hdo_kravspec} uten ekstra vedlikeholdsarbeid.

For observability ble \texttt{prometheus-net} valgt fordi det integrerer
tett med ASP.NET Core og støtter eksponering av både innebygde og
egendefinerte metrikker til Prometheus, i tråd med krav O1.3
\cite{hdo_kravspec}. Docker og Docker Compose ble brukt for
containerisering i henhold til krav I1.1 \cite{hdo_kravspec}. GitHub
Actions ble valgt for CI/CD fordi kildekoden allerede lå på GitHub og
verktøyet ikke krever ekstern infrastruktur.

Tabell~\ref{tab:avhengigheter} gir en oversikt over de sentrale
pakkeavhengighetene som ble benyttet.

\begin{table}[H]
\centering
\caption{Sentrale pakkeavhengigheter i HDO-VideoAPI.}
\label{tab:avhengigheter}
\begin{tabular}{llr}
\toprule
\textbf{Pakke} & \textbf{Formål} & \textbf{Versjon} \\
\midrule
\multicolumn{3}{l}{\textit{Hovedprosjekt}} \\
\texttt{Microsoft.AspNetCore.Authentication.JwtBearer} & JWT-autentisering   & \texttt{10.0.5}  \\
\texttt{Microsoft.IdentityModel.Tokens}                & JWT-validering      & \texttt{8.16.0}  \\
\texttt{System.IdentityModel.Tokens.Jwt} & JWT-håndtering   & \texttt{8.16.0}  \\
\texttt{Microsoft.AspNetCore.OpenApi} & OpenAPI   & \texttt{10.0.2}  \\
\texttt{Swashbuckle.AspNetCore}                        & OpenAPI/Swagger     & \texttt{10.1.1}  \\
\texttt{prometheus-net.AspNetCore}                     & HTTP-metrikk        & \texttt{8.2.1}   \\
\texttt{prometheus-net.AspNetCore.HealthChecks}        & Helsesjekk-metrikk  & \texttt{8.2.1}   \\
\texttt{DotNetEnv}                                     & Lokal .env-støtte   & \texttt{3.1.1}   \\
\midrule
\multicolumn{3}{l}{\textit{Testprosjekt}} \\
\texttt{xunit}                                         & Testrammeverk       & \texttt{2.9.3}   \\
\texttt{Moq}                                           & Mocking             & \texttt{4.20.72} \\
\texttt{Microsoft.AspNetCore.Mvc.Testing}              & Integrasjonstesting & \texttt{10.0.3}  \\
\texttt{coverlet.collector}                            & Kodedekning         & \texttt{6.0.4}   \\
\bottomrule
\end{tabular}
\end{table}